% Generated by `taskgen course examples` from ecc.double_and_add_ternary (seed 1).
% Do not edit: change examples.yaml and regenerate.
\begin{example}[тот же показатель, тернарное разложение]\label{ex:ecdlp-ternary}
Вычислить степень $47P$ точки $P = (1, 52)$ эллиптической кривой \[y^2 \equiv x^3 + 50x + 80\ (\textrm{mod}\ 83)\] с помощью тернарного разложения показателя $47$.

\par\noindent\textit{Решение.} Тернарное разложение (цифры $-1, 0, 1$, без соседних ненулевых): $47 = 10\bar 1000\bar 1$ ($\bar 1 = -1$), т.е. $47 = 2^{6} - 2^{4} - 2^{0}$; для цифры $-1$ прибавляем $-P = (1, 31)$. Начиная с $R = P$, для каждой следующей цифры: $R \gets 2R$, затем при ненулевой цифре $R \gets R \pm P$. Формулы: $\lambda = \frac{3x_1^2 + a}{2y_1}$ (удвоение), $\lambda = \frac{y_2 - y_1}{x_2 - x_1}$ (сложение), $x_3 = \lambda^2 - x_1 - x_2$, $y_3 = \lambda(x_1 - x_3) - y_1$; $a = 50$, всё по модулю $83$.\[\begin{array}{l|r|r|c}\text{операция} & kP & \lambda & R \\ \hline 2R & 2P & 46 & (39, 26) \\ 2R & 4P & 52 & (53, 76) \\ R - P & 3P & 28 & (66, 58) \\ 2R & 6P & 63 & (19, 81) \\ 2R & 12P & 28 & (82, 64) \\ 2R & 24P & 27 & (67, 9) \\ 2R & 48P & 27 & (14, 11) \\ R - P & 47P & 24 & (63, 58) \\ \end{array}\]\[47P = (63, 58).\]
\end{example}
